% --------------------------------- Preamble -----------------------------------
\documentclass[11pt, letterpaper, twocolumn]{article}

\usepackage[T1]{fontenc} % use 8-bit fonts
\usepackage[margin=0.75in, columnsep=0.2in, % 3.4 in columns
    headsep=0.28in, footskip=0.42in]{geometry} % paper sizes
\usepackage{amsmath, amsfonts, amssymb} % improved math
\usepackage{bm} % bold math (must be after amsmath)
\usepackage{upquote} % upright single quotes in verbatim
\usepackage{pgfplots} % includes tikz, xcolor, graphicx
\usepackage[colorlinks=true, allcolors=gray]{hyperref}

\raggedbottom % do not force stretching text to fill the page
\setlength{\emergencystretch}{\hsize} % avoid overfull lines
\overfullrule=1mm % black band to show overfull lines
\pdfsuppresswarningpagegroup=1 % remove pdf image warnings
\frenchspacing % consistent spacing between words
\newcommand{\ul}[1] % vectors
    {{}\mkern1mu\underline{\mkern-1mu#1\mkern-1mu}\mkern1mu}
\DeclareSymbolFont{euler}{U}{eur}{m}{n} % symbol font
\DeclareMathSymbol{\PI}{\mathord}{euler}{25} % up pi
\pgfplotsset{compat=newest} % use newest pgfplot version
\newcommand{\EE}[2]{\ensuremath{{{#1}\!\times\!10^{#2}}}}
% ------------------------------------------------------------------------------

\title{Extended Kalman Filter for Orbital State Estimation}
\author{}
\date{} 

\begin{document}
\maketitle

\section*{Description}

This project implements an Extended Kalman Filter for estimating the state of a 
satellite in 2D orbit about the Earth. The system dynamics are nonlinear due to
gravitational acceleration, and the measurement model is nonlinear because the 
sensor provides only the bearing angle to the satellite. The filter estimates
position and velocity over an eight-hour orbit using recursive linearization
of both the dynamics and measurement models.

\subsection*{Problem Formulation}

A satellite orbits the Earth. For simplicity, we will say that the gravitational acceleration 
on the satellite is $G_E/r^2$, where we use the value $G_E=3.986\times10^{14} m^3/s^2$ and 
$r$ is the total distance from earth to the satellite, $r=\sqrt{p_x^2+p_y^2}$. The earth is at 
position (0,0) and at this origin is a sensor that only measures the bearing to the satellite. 
Our state vector is the $x$ and $y$ position and velocity values:
\[
\ul{x} = \begin{bmatrix}
    p_x & p_y & v_x & v_y
\end{bmatrix}^T
\]
The acceleration of $G_E/r^2$ is just the magnitude. The vector of acceleration is the 
magnitude multiplied by the negative unit vector of position: $p_x/r$ and $p_y/r$:
\begin{figure}[h]
    \includegraphics[width=\linewidth]{../figures/fig_force.pdf}
\end{figure}

So, the dynamics equations are:
\begin{center}
    $\dot{p_x} = v_x + \nu_{p_x}$
    \\
    $\dot{p_y} = v_y + \nu_{p_y}$
    \\
    $\dot{v_x} = a_x + \nu_{v_x}$
    \\
    $\dot{v_y} = a_y + \nu_{v_y}$
\end{center}
where $a_x$ and $a_y$ are defined as
\[\begin{aligned}
a_x &= -\frac{G_Ep_x}{(p_x^2+p_y^2)^{3/2}} \\
a_y &= -\frac{G_Ep_y}{(p_x^2+p_y^2)^{3/2}}
\end{aligned}
\]
and where the $\ul{\nu} \sim N(0,\mathbf{Q})$ values are noise. In the discrete domain, using
forward Euler integration, the dynamics equations are:
\begin{center}
    $p_{x,k+1} = p_{x,k} + v_{x,k}T + \nu_{p_x}\sqrt{T}$
    \\
    $p_{y,k+1} = p_{y,k} + v_{y,k}T + \nu_{p_y}\sqrt{T}$
    \\
    $v_{x,k+1} = v_{x,k} + a_{x,k}T + \nu_{v_x}\sqrt{T}$
    \\
    $v_{y,k+1} = v_{y,k} + a_{y,k}T + \nu_{a_y}\sqrt{T}$
\end{center}
The measurement equation is:
\begin{center}
    $z = \operatorname{atan2}(p_x,p_y) + \eta$
\end{center}
where $\eta \sim N(0,R)$.

\section*{Extended Kalman Filter}

Because both the dynamics and measurement models are nonlinear, state estimation is performed
with an Extended Kalman Filter. At each time step, the current estimate is propagated 
through the nonlinear dynamics, and the covariance is propagated using the Jacobian of 
the dynamics.

The measurement update uses the nonlinear bearing model together with the Jacobian of the
measurement function. The innovation is computed as the difference between the actual
and predicted bearing measurements. Since bearing is an angle, the innovation must be
wrapped to the interval $[-\pi,\pi)$ before applying the update.

\subsection*{Dynamics Jacobian}

The nonlinear dynamics are
\[
\phi(x) = \begin{bmatrix}
    p_x + v_xT \\
    p_y + v_yT \\
    v_x + a_xT \\
    v_y + a_yT
\end{bmatrix}
\]

The Jacobian of the dynamics is
\[ \bm{\varPhi} = \begin{bmatrix}
    1 & 0 & T & 0 \\
    0 & 1 & 0 & T \\
    G_ET\frac{2p_x^2-p_y^2}{(p_x^2+p_y^2)^{5/2}} & \frac{3G_ETp_xp_y}{(p_x^2+p_y^2)^{5/2}} & 1 & 0 \\
    \frac{3G_ETp_xp_y}{(p_x^2+p_y^2)^{5/2}} & G_ET\frac{2p_y^2-p_x^2}{(p_x^2+p_y^2)^{5/2}} & 0 & 1
\end{bmatrix} 
\]

\subsection*{Measurement Jacobian}

The measurement model is
\[
h(x) = \operatorname{atan2}(p_y,p_x)
\]

The Jacobian of the measurement model is
\[ \bm{H} = \begin{bmatrix}
    \frac{-p_y}{p_x^2+p_y^2} & \frac{p_x}{p_x^2+p_y^2} & 0 & 0
\end{bmatrix} 
\]

\section*{Results}

Figure~\ref{fig:path} shows the true orbital path together with the estimated path
from the Extended Kalman Filter.

\begin{figure}[h]
    \centering
    \includegraphics[width=\linewidth]{../figures/fig_path.pdf}
    \caption{True and estimated orbital paths}
    \label{fig:path}
\end{figure}

\pagebreak

Figures~\ref{fig:poserr} and \ref{fig:velerr} show the position and velocity estimation
errors over time along with their corresponding $3-\sigma$ bounds from the filter covariance.

\begin{figure}[h]
    \centering
    \includegraphics[width=\linewidth]{../figures/fig_pos_errors.pdf}
    \caption{Position estimation errors}
    \label{fig:poserr}
\end{figure}

\begin{figure}[h]
    \centering
    \includegraphics[width=\linewidth]{../figures/fig_vel_errors.pdf}
    \caption{Velocity estimation errors}
    \label{fig:velerr}
\end{figure}

\section*{Discussion}

The results show that the Extended Kalman Filter is able to track the satellite
state despite nonlinear dynamics and a nonlinear bearing-only measurement model.
The position and velocity errors remain bounded, and comparison with the $3-\sigma$
bounds provides insight into the consistency of the covariance propagation.

A key part of this implementation is the use of analytical Jacobians for both the
dynamics and measurement models. Another important detail is angle wrapping in the
innovation, since unwrapped bearing residuals can create large artificial errors 
and destabilize the filter.

This project highlights the central EKF idea: a nonlinear estimation problem can be
handled by propagating the full nonlinear state while locally linearizing the system
for covariance calculations.

\end{document}